\documentclass[11pt]{article}
\usepackage[margin=1in]{geometry}
\usepackage{booktabs}
\usepackage{array}
\usepackage{enumitem}
\usepackage[hidelinks]{hyperref}
\setlength{\parskip}{0.5em}
\setlength{\parindent}{0pt}
\renewcommand{\arraystretch}{1.15}

\title{Progress Report -- Week [4]}
\author{%
  \begin{tabular}[t]{c@{\hspace{2em}}c@{\hspace{2em}}c}
    Aniyah Mack & Adonay Zeru & Hasset Admasu\\
    \texttt{amack18@emich.edu} & \texttt{azeru@emich.edu} & \texttt{hadmasu@emich.edu}
  \end{tabular}\\[0.8em]
  Nathnael Abraham\\
  \texttt{abrah10@emich.edu}}
\date{October 12, 2026}

\begin{document}
\maketitle

\section*{Project Overview}
Our project investigates how vision-language models can identify important events in soccer match
videos and generate natural-language commentary describing those events. Our main goal is to
determine how accurately soccer events can be identified from video, and how those events can be
used to generate accurate commentary for different audiences, including casual viewers, expert fans,
and people with low vision. The project builds on the SoccerReplay-1988 dataset~[5] and existing
baseline methods for soccer event classification.

\section*{Where Our Last Report Left Off}
Last week we reported no results on real data. We had built a supervised evaluation (game-level split,
cleaned labels, several seeds, macro-F1) and tested it on synthetic features, while the videos and
features were still being re-obtained. We had set the contrastive and temporal work aside, on the
advice to make a mean-pooled baseline solid first.

\section*{Current Progress}
\subsection*{Cleanup and moving our compute to the professor's server}
Most of this period went into cleaning up the project and moving all heavy computation off our personal
machines onto the professor's GPU server.
\begin{itemize}[noitemsep]
  \item \textbf{Cleanup.} We reorganized the \texttt{dino-contrastive} code into a conventional layout
        (a shared package, one script per stage, one entry point), removed superseded scripts, took
        credentials out of the code, pinned the dependencies, and added automated tests and a secret scan.
        The cleaned code lives on its own branch, \texttt{cleanup/dino-contrastive}, kept locally and not yet
        pushed. Every removal was listed and approved before it was made.
  \item \textbf{Moving compute to the server.} Match videos are now downloaded straight from the shared
        folder to the server, frames are extracted there and the videos deleted, and encoding, clustering and
        training all run there. Nothing heavy runs on our own computers any more. The same command
        reproduces a run from a fresh copy of the code, and we checked that it gives identical frames.
  \item \textbf{Audit.} The cleanup exposed problems in the older pipeline (listed below), which we fixed
        before producing the results that follow.
\end{itemize}

\subsection*{Results}
Once the pipeline was running on the server we produced real results, from the same
\texttt{dino-contrastive} pipeline rather than the separate baseline repository of last week. Every number
below is produced by it.

\textbf{Data.} 65 English Premier League 2021--22 matches, split by whole match as in SoccerReplay's
official split (50 train, 7 validation, 8 test). Each event is described by up to 60 frames (30\,s before to
30\,s after, one per second), encoded with frozen DINOv2-base. After removing junk and non-visual
labels there are 4{,}337 events of 22 types (3{,}329 train, 472 validation, 536 test). We did not drop
rare types: the smallest are \texttt{second\_yellow\_card} (a single event in the whole dataset, so it can
never be scored), \texttt{throw\_in} (7) and \texttt{penalty} (12). Corners are the largest type, 16.4\% of
events.

\textbf{Protocol.} Models are fitted on training matches, chosen on validation matches, and scored once
on test matches. Because 8 test matches are few, we also ran \emph{5-fold cross-validation by match}:
every one of the 65 matches is a test match exactly once (13 per fold, 7 validation matches, 45 training
matches), with 2 seeds per fold. Intervals come from resampling whole matches.

\subsection*{Event classification}
Table~\ref{tab:cv} is our main result. Table~\ref{tab:official} gives the official split for comparison.
The sequence model is the temporal model chosen on validation among four candidates; it was the
bidirectional LSTM in every fold and in the official split.

\begin{table}[h]
\centering\small
\begin{tabular}{lcccc}
\toprule
Model & Accuracy & Macro-F1 & All 65 matches pooled (95\% interval) \\
\midrule
Always predict the most common label & $0.164 \pm 0.012$ & -- & -- \\
Pooled frozen features & $0.371 \pm 0.019$ & $0.239 \pm 0.014$ & $0.371$ ($0.358$--$0.384$) \\
Sequence model (BiLSTM) & $\mathbf{0.463 \pm 0.014}$ & $\mathbf{0.318 \pm 0.014}$ & $0.462$ ($0.447$--$0.479$) \\
Visual + annotation features & $0.382 \pm 0.012$ & $0.269 \pm 0.017$ & $0.382$ ($0.368$--$0.398$) \\
\bottomrule
\end{tabular}
\caption{5-fold cross-validation by match (mean $\pm$ sd over folds). The sequence model beats the pooled
features by $0.092 \pm 0.011$ accuracy, and is ahead in all 5 folds.}
\label{tab:cv}
\end{table}

\begin{table}[h]
\centering\small
\begin{tabular}{lccc}
\toprule
Model & Accuracy & Macro-F1 & 95\% interval (accuracy)\\
\midrule
Always predict the most common label & $0.166$ & -- & --\\
Pooled frozen features & $0.390 \pm 0.020$ & $0.247 \pm 0.017$ & $0.353$--$0.428$\\
Sequence model (BiLSTM) & $0.491 \pm 0.022$ & $0.324 \pm 0.032$ & $0.464$--$0.522$\\
Visual + annotation features & $0.421 \pm 0.013$ & $0.296 \pm 0.041$ & $0.382$--$0.461$\\
\bottomrule
\end{tabular}
\caption{Official split, 8 test matches (536 events); mean $\pm$ sd over 5 seeds. The sequence model is ahead of
the pooled features in all 5 seeds.}
\label{tab:official}
\end{table}

Accuracy flatters the common types, so macro-F1 is the fairer number: even the best model reaches only
about 0.32. Table~\ref{tab:perclass} shows why. Seven of the 21 scorable types are essentially never recognized
(F1 below 0.10), and they are the rare ones; start and end of a half, goals and substitutions are recognized
well.

\begin{table}[h]
\centering\footnotesize
\begin{tabular}{lrcccc}
\toprule
Event type & Events & Precision & Recall & F1 & \\
\midrule
corner & 712 & 0.52 & 0.66 & 0.58 & \\
clearance & 631 & 0.39 & 0.51 & 0.44 & \\
shot\_off\_target & 380 & 0.31 & 0.39 & 0.35 & \\
foul\_with\_no\_card & 367 & 0.35 & 0.39 & 0.37 & \\
substitution & 355 & 0.68 & 0.66 & 0.67 & \\
goal & 266 & 0.65 & 0.66 & 0.65 & \\
lead\_to\_corner & 264 & 0.42 & 0.32 & 0.36 & \\
yellow\_card & 224 & 0.38 & 0.39 & 0.38 & \\
saved\_by\_goal-keeper & 192 & 0.18 & 0.05 & 0.08 & low\\
injury & 157 & 0.49 & 0.43 & 0.45 & \\
ball\_possession & 151 & 0.39 & 0.17 & 0.24 & \\
end\_of\_half\_game & 131 & 0.67 & 0.66 & 0.66 & \\
start\_of\_half\_game & 130 & 0.78 & 0.87 & 0.82 & \\
off\_side & 119 & 0.21 & 0.11 & 0.15 & \\
var & 71 & 0.37 & 0.22 & 0.28 & \\
free\_kick & 66 & 0.19 & 0.06 & 0.09 & low\\
ball\_out\_of\_play & 53 & 0.00 & 0.00 & 0.00 & low\\
foul\_lead\_to\_penalty & 30 & 0.17 & 0.05 & 0.08 & low\\
red\_card & 18 & 0.00 & 0.00 & 0.00 & low\\
penalty & 12 & 0.00 & 0.00 & 0.00 & low\\
throw\_in & 7 & 0.00 & 0.00 & 0.00 & low\\
\bottomrule
\end{tabular}
\caption{Per-class results of the sequence model, pooling the test predictions of all 65 matches from the
cross-validation (both seeds; event counts are per seed). The 22nd type, \texttt{second\_yellow\_card}, has a
single event and is left out.}
\label{tab:perclass}
\end{table}

\subsection*{Do the embeddings follow the match or the event?}
This was our main finding last time, and it holds at 65 matches. For K-means with one cluster per event
type, raw mean-pooled DINOv2 embeddings agree with the match about as much as with the event (ARI
0.069 against 0.036). On matches never seen in training the gap is much larger: match ARI 0.354 against
event ARI 0.035. Normalizing each match separately removes the match signal (match ARI 0.003, event ARI
0.046). Clustering inside each match gives purity 0.446 against 0.374 for shuffled clusters.

\subsection*{Contrastive head}
A triplet-loss projection head (256 dimensions), trained on training matches only, removes the match signal
on held-out matches (match ARI $0.354 \to 0.029$) and raises event agreement (event ARI $0.035 \to 0.106$,
purity $0.278 \to 0.370$; up in all 5 seeds). In cross-validation the event ARI goes from $0.031 \pm 0.002$ to
$0.094 \pm 0.003$, up in all 5 folds. The improvement is real but small in absolute terms. On the training
matches themselves the event ARI reaches 0.73, so the head mostly memorizes there.

\subsection*{What went wrong, and what we fixed}
Most of this period went into making these numbers trustworthy. We found and fixed, each with a test that
fails on the old code:
\begin{itemize}[noitemsep]
  \item Our earlier contrastive result (96\% purity) was invalid: the head had been trained using test-game
        labels. Training and evaluation now use disjoint matches, and the training code refuses overlap.
  \item About 47\% of events had been cut from the wrong half's video, because match timestamps restart in
        each half. Fixed.
  \item Some videos report a wrong frame rate, so counting frames drifted away from the video's clock
        (up to 21 minutes by the end of a half). Frames are now chosen by timestamp, and 15.6\% of all frames
        (42 of 130 halves) changed. Rerunning moved every result by no more than seed noise.
  \item Smaller issues: annotation features that could see later events, a BiLSTM reading padding, and
        evaluation choices made on test data. All 22 findings are logged with their effect.
\end{itemize}

\section*{Response to Feedback}
\begin{table}[h]
\centering\small
\begin{tabular}{>{\raggedright\arraybackslash}p{4.0cm}>{\raggedright\arraybackslash}p{2.3cm}>{\raggedright\arraybackslash}p{8.2cm}}
\toprule
Feedback & Status & What we did\\
\midrule
1. Make the mean-pooled pipeline and evaluation solid first & Done & Match-level splits, cross-validation, validation-only model choice, several seeds, intervals from resampling matches, and the mean-pooled baseline on real features (Tables~\ref{tab:cv}, \ref{tab:official}).\\
2. Sweep embedding size (64--512) on held-out data & Not started & The head is still fixed at 256 dimensions.\\
3. Address imbalance before hard-negative mining & Not started & Anchors still follow class frequency. We now know which types fail (Table~\ref{tab:perclass}), which makes this the next step.\\
4. Hold out whole games, ideally a season or competition & Partly done & Whole matches are held out everywhere, and we also ran a leave-stadium-out split (no test stadium in training; results close to the match split). A second season needs new videos.\\
5. Report supervised metrics, especially per class & Done & Macro-F1, per-class precision, recall and F1 (Table~\ref{tab:perclass}), and confusion matrices for every model.\\
6. Do not move to LSTM/Transformer yet & \textbf{Not followed} & We did run BiLSTM, temporal-CNN and Transformer models on the same split. They beat the pooled baseline by about 9--10 points, but this goes against your advice, so we present them only as a comparison and ask below whether to keep them.\\
\bottomrule
\end{tabular}
\caption{Feedback points and our response.}
\end{table}

\section*{Challenges and Issues}
\textbf{Hidden data problems took most of the time.} The bugs above each looked like a plausible modelling
result until we audited the data. We now check frame timing on every video and state it in every report.

\textbf{Compute.} The shared GPU is usually full, so our runs use the CPU: encoding the 65 matches takes
hours, and the cross-validation took about 3 hours.

\textbf{Few test matches.} Eight test matches leave real uncertainty, which is why cross-validation is our main
number. Seeds alone understate it.

\textbf{Data access.} Dropbox access tokens expired every 4 hours; the downloads now renew themselves.

\section*{Next Steps}
\begin{enumerate}[noitemsep]
  \item Sweep the embedding size (64, 128, 256, 512), compared on held-out macro-F1 in cross-validation.
  \item Add class-balanced anchor sampling, then compare random and semi-hard negatives by per-class F1,
        aiming at the seven rare types that are never recognized.
  \item Decide how to treat very rare types (merge or keep), based on your answer below.
  \item Hold out a second league-season or a different competition, which needs new videos and features.
  \item Only then scale beyond 65 matches, and revisit temporal models against the final baseline.
\end{enumerate}

\section*{Questions or Feedback Needed}
\begin{itemize}[noitemsep]
  \item We ran sequence models despite the advice to wait. Should we keep them as a comparison against the
        pooled baseline, or drop them until the contrastive steps are finished?
  \item Is 5-fold cross-validation by match over 65 matches an acceptable main result, with the official
        split as secondary?
  \item Seven rare event types are almost never recognized. Should they be merged into broader categories
        (for example red card into card), or kept as they are?
  \item For the held-out competition, is another EPL season enough, or should it be a different competition?
\end{itemize}

\section*{Team Contributions}
\begin{itemize}[noitemsep]
  \item \textbf{Aniyah Mack:} [to be filled in by Aniyah]
  \item \textbf{Adonay Zeru:} [to be filled in by Adonay]
  \item \textbf{Hasset Admasu:} Built and audited the \texttt{dino-contrastive} pipeline: leakage and data
        fixes (wrong half, frame timing, annotation features), match-level splits and 5-fold
        cross-validation, the sequence and contrastive experiments, a GPU-server workflow with tests and
        continuous integration, and the analysis and results in this report.
  \item \textbf{Nathnael Abraham:} [to be filled in by Nathnael]
\end{itemize}

\section*{References}
\begin{enumerate}[label={[\arabic*]}, noitemsep]
  \item Claude Code. Helped with identifying the data leakage, pipeline implementation, analysis and documentation.
  \item James MacQueen. Some methods for classification and analysis of multivariate observations. In
        \emph{Proceedings of the Fifth Berkeley Symposium on Mathematical Statistics and Probability, Volume 1:
        Statistics}, pages 281--297. University of California Press, 1967.
  \item Maxime Oquab, Timoth\'ee Darcet, Th\'eo Moutakanni, Huy V. Vo, Marc Szafraniec, Vasil Khalidov, Pierre
        Fernandez, Daniel Haziza, Francisco Massa, Alaaeldin El-Nouby, Mahmoud Assran, Nicolas Ballas, Wojciech
        Galuba, Russell Howes, Po-Yao Huang, Shang-Wen Li, Ishan Misra, Michael Rabbat, Vasu Sharma, Gabriel
        Synnaeve, Hu Xu, Herv\'e J\'egou, Julien Mairal, Patrick Labatut, Armand Joulin, and Piotr Bojanowski.
        DINOv2: Learning robust visual features without supervision. \emph{Transactions on Machine Learning
        Research}, 2024.
  \item Alec Radford, Jong Wook Kim, Chris Hallacy, Aditya Ramesh, Gabriel Goh, Sandhini Agarwal, Girish Sastry,
        Amanda Askell, Pamela Mishkin, Jack Clark, Gretchen Krueger, and Ilya Sutskever. Learning transferable
        visual models from natural language supervision. In \emph{Proceedings of the 38th International
        Conference on Machine Learning (ICML)}, volume 139, pages 8748--8763. PMLR, 2021.
  \item Jiayuan Rao, Haoning Wu, Hao Jiang, Ya Zhang, Yanfeng Wang, and Weidi Xie. Towards universal soccer
        video understanding. In \emph{Proceedings of the IEEE/CVF Conference on Computer Vision and Pattern
        Recognition (CVPR)}, pages 8384--8394, 2025.
  \item Florian Schroff, Dmitry Kalenichenko, and James Philbin. FaceNet: A unified embedding for face
        recognition and clustering. In \emph{Proceedings of the IEEE Conference on Computer Vision and Pattern
        Recognition (CVPR)}, pages 815--823, 2015.
  \item Aniyah Mack, Adonay Zeru, Hasset Admasu, and Nathnael Abraham. Soccer Event Baseline: supervised
        mean-pooled baseline for soccer event recognition. GitHub repository,
        \url{https://github.com/Hasset-hub/soccer-event-baseline}, 2026.
  \item Laurens van der Maaten and Geoffrey Hinton. Visualizing data using t-SNE. \emph{Journal of Machine
        Learning Research}, 9(86):2579--2605, 2008.
\end{enumerate}

\end{document}
